% Generated by scripts/assemble.py; corrections belong in source/corrections.json.
\examhead{104}{中国科学技术大学 2024 年数学分析试卷}{103}
% p104-b05; source PDF page 104
\blockanchor{p104-b05}
\question{1}(15 分) 设 $f = f(u, v, w)$ 是 $\mathbb{R}^3$ 上定义的光滑函数, 求 $g(x, y, z) = f\left(x + y, xy, \frac{x}{y}\right)$ 的偏导数 $\frac{\partial g}{\partial x}, \frac{\partial^2 g}{\partial x^2}$.\par
% p104-b06; source PDF page 104
\blockanchor{p104-b06}
\question{2}(15 分) 设 $a > 0$, 求旋轮线一拱 $\left\{ \begin{array}{l} x = a(t - \sin t); \\ y = a(1 - \cos t). \end{array} \right.$$0 \leq t \leq 2\pi$ 的弧长.\par
% p104-b07; source PDF page 104
\par\Needspace{6\baselineskip}
\blockanchor{p104-b07}
\question{3}(15 分) 计算积分\par
\nopagebreak[4]
% p104-b08; source PDF page 104
\blockanchor{p104-b08}
\displayformula{\iint_ {S} (x ^ {2} + y) \mathrm{d} y \mathrm{d} z - y \mathrm{d} z \mathrm{d} x + (x ^ {2} + y ^ {2}) \mathrm{d} x \mathrm{d} y.}
% p104-b09; source PDF page 104
\blockanchor{p104-b09}
其中 $S$ 为旋转抛物面 $z = x^{2} + y^{2}$ 在 $z \leq 1$ 的部分, 并取曲面外侧为正向.\par
% p104-b10; source PDF page 104
\blockanchor{p104-b10}
\question{4}(15 分) 计算极限  $\lim_{n\to\infty} n\left(\frac{1}{n^{2}+1^{2}}+\frac{1}{n^{2}+2^{2}}+\cdots+\frac{1}{n^{2}+n^{2}}\right)$ .\par
% p104-b11; source PDF page 104
\blockanchor{p104-b11}
\question{5}(15 分) 确定函数 $I(\alpha) = \int_{0}^{+\infty} \frac{\ln(1 + x^3)}{x^\alpha} \, \mathrm{d}x$ 的定义域, 并证明 $I$ 在定义域上是连续函数.\par
% p104-b12; source PDF page 104
\blockanchor{p104-b12}
\question{6}(15 分) 设 $a > 0$, $x_0 > 0$ 而 $x_n$ 由递推公式 $x_n = \frac{1}{2} \left( x_{n-1} + \frac{a}{x_{n-1}} \right)$ ($n \geq 1$) 给出. 证明数列 $\{x_n\}_{n \geq 1}$ 是单调数列, 并求 $\lim_{n \to \infty} x_n$.\par
% p104-b13; source PDF page 104
\blockanchor{p104-b13}
\question{7}(15 分) 设 $f: \mathbb{R}^2 \to \mathbb{R}$ 在 $(0,0)$ 附近存在偏导数 $\frac{\partial f}{\partial x}$ 和 $\frac{\partial f}{\partial y}$. 设 $\frac{\partial f}{\partial x}$ 在 $(0,0)$ 点是连续的. 求证: 对于任何 \displayformula{u, v \in \mathbb{R}, \left.\frac{\mathrm{d}}{\mathrm{d}t}\right|_{t=0} f(tu, tv) = u \frac{\partial f}{\partial x}(0,0) + v \frac{\partial f}{\partial y}(0,0)}.\par
% p104-b14; source PDF page 104
\blockanchor{p104-b14}
\question{8}(15 分) 设 $f(x)$ 在 $[-1, 1]$ 上二次可导, 并有 $f(-1) = 0$, $f(0) = 0$, $f(1) = 1$. 证明: 存在 $\xi \in (-1, 1)$, 使得 $f''(\xi) = 1$.\par
% p104-b15; source PDF page 104
\par\Needspace{6\baselineskip}
\blockanchor{p104-b15}
\question{9}(15 分) 设 $f(x)$ 是 $2\pi$ 为周期的黎曼可积函数, 设 $a_{n}, b_{n}$ 是它的 Fourier 系数. 求证: 对于任何的 $c, d \in [-\pi, \pi]$, 都有\par
\nopagebreak[4]
% p104-b16; source PDF page 104
\blockanchor{p104-b16}
\displayformula{\int_ {c} ^ {d} f (x) \mathrm{d} x = \int_ {c} ^ {d} \frac {a _ {0}}{2} \mathrm{d} t + \sum_ {n = 1} ^ {\infty} \int_ {c} ^ {d} \left(a _ {n} \cos n x + b _ {n} \sin n x\right) \mathrm{d} x.}
% p104-b17; source PDF page 104
\par\Needspace{6\baselineskip}
\blockanchor{p104-b17}
\question{10}(15 分) 设有实数列 $\{a_{n}\}_{n\geq 1}$, 令\par
\nopagebreak[4]
% p104-b18; source PDF page 104
\blockanchor{p104-b18}
\displayformula{S _ {n} = a _ {1} + a _ {2} + \dots + a _ {n}, \sigma_ {n} = \frac {S _ {1} + S _ {2} + \cdots + S _ {n}}{n}.}
% p104-b19; source PDF page 104
\blockanchor{p104-b19}
(1) 证明: $\{S_n\}$ 有极限推出 $\{\sigma_n\}$ 有相同的极限.\par
% p104-b20; source PDF page 104
\blockanchor{p104-b20}
(2) 如果 $\{\sigma_n\}$ 收敛, 而且 $a_n = o\left(\frac{1}{n}\right)$, 则 $\{S_n\}$ 收敛.\par
